\chapter{Marco teorico}

\subsection{Gemelos digitales en sistemas electricos de potencia}

Un gemelo digital es una representacion computacional sincronizada con un
activo fisico que permite reproducir su comportamiento en tiempo real y
explorar escenarios hipoteticos~\citep{Tao2018,TwinPower2020}. En sistemas
electricos de potencia, los gemelos digitales aparecieron primero en
transmision para estudios de flujos de potencia y estabilidad transitoria;
mas recientemente la literatura se ha enfocado en subestaciones de
distribucion~\citep{DT-Substation-2022} y en plantas renovables.

Desde el punto de vista formal, un gemelo digital para una subestacion
puede verse como un mapeo

\[
\mathcal{T}: (\theta, u_t) \mapsto \hat{x}_t
\]

donde $\theta$ son los parametros de la subestacion (topologia, parametros
de transformadores y lineas), $u_t$ son las consignas operativas (tap,
referencia de tension, despacho) y $\hat{x}_t$ es la prediccion del estado
electrico esperado al tiempo $t$. La funcion $\mathcal{T}$ se obtiene
combinando modelos de estado estacionario (flujo de potencia), modelos
dinamicos (swing de orden reducido) y modelos termicos (primer orden con
constante de tiempo de 60-90~min).

En este trabajo se adopta un gemelo digital \emph{ligero} (modelo de bajo
orden, sin componentes de alta frecuencia) por tres motivos:

\begin{enumerate}[leftmargin=*,itemsep=2pt]
  \item \textbf{Reducir el costo de calibracion}. Un gemelo detallado
        (incluyendo armonicos de conmutacion, modelos de bujes y
        modelos termicos distribuidos) requiere meses de datos
        historicos para identificar parametros, mientras que un gemelo
        ligero se identifica con la hoja de datos del fabricante y
        unas pocas semanas de operacion normal.

  \item \textbf{Generar datos sinteticos}. El gemelo ligero es
        suficientemente rico para que las firmas de los ocho tipos de
        falla modelados en el capitulo 3 sean distinguibles en los
        residuos.

  \item \textbf{Computo en linea}. La prediccion del gemelo se ejecuta
        a cadencia SCADA (1~s), lo que exige un modelo con pocos
        parametros y baja latencia.
\end{enumerate}

\subsection{Conformal prediction}

La \emph{conformal prediction}~\citep{Vovk2005,Shafer2008} es un marco
para construir regiones de prediccion con garantia de cobertura bajo
el supuesto debil de intercambiabilidad (no exige distribucion
conhecida). En el contexto de deteccion de anomalias, se aplica como
sigue.

Sea $e_1, \dots, e_n$ una secuencia de reconstruction errors producida
por el autoencoder sobre datos de calibracion. Bajo intercambiabilidad,
la fraccion de datos de test con reconstruction error por encima del
cuantil empirico $q_{1-\alpha} = \text{Quantil}_{n+1}^{\lceil (1-\alpha)(n+1) \rceil}$
es a lo mas $\alpha$. Es decir,

\[
\mathbb{P}\{e_{n+1} > q_{1-\alpha}\} \le \alpha.
\]

Esta garantia es fundamental para que el sistema sea aceptado por el
operador: las falsas alarmas no se acumular bajo cambios de regimen
operativo, mientras que las fallas reales siguen detectandose con alta
probabilidad.

\subsubsection{Por que conformal y no solo un umbral 3-sigma}

Un umbral 3-sigma asume gaussianidad, lo que es violado en la practica
por los picos de carga, los cambios de tap y los eventos climaticos.
La consecuencia es una tasa de falsa alarma muy superior al objetivo
(el detector 3-sigma empirico del capitulo 5 alcanza un FPR de 13.2\%).
Conformal entrega una garantia libre de distribucion y, en este
trabajo, se demuestra empiricamente que cumple el objetivo.

\subsection{Aprendizaje federado}

El aprendizaje federado~\citep{McMahan2017} entrena un modelo compartido
entre multiples clientes sin mover los datos crudos. FedAvg, el
algoritmo mas usado, realiza los siguientes pasos en cada ronda $t$:
\enumerate{
  \item El servidor envia el modelo global $\theta_t$ a cada cliente $k$.
  \item Cada cliente entrena $\theta_t$ sobre sus datos locales $D_k$
        durante $E$ epocas, obteniendo $\theta^k_{t+1}$.
  \item El servidor agrega los modelos locales:

\[
\theta_{t+1} = \sum_{k=1}^{K} \frac{n_k}{N} \theta^k_{t+1}
\]
}
donde $n_k = |D_k|$ y $N = \sum_k n_k$.

En este caso la federacion FedAvg entre tres subestaciones (prototipo;
en produccion se usan cinco o mas) tiene dos
beneficios:
\begin{itemize}[leftmargin=*,itemsep=2pt]
  \item \textbf{Estadistica}: el modelo global ve la variabilidad
        geografica (perfiles de carga, clima, mix de clientes) que
        ninguna SE ve por si sola.
  \item \textbf{Cumplimiento}: los datos nunca abandonan la SE, lo
        que cumple con la guia de la SEC para redes OT y con NERC
        CIP (que las distribuidoras chilenas usan como benchmark
        internacional).
\end{itemize}

\subsubsection{Limitaciones}

FedAvg asume clientes \emph{iid}. En distribucion, las SE tienen
heterogeneidad fuerte (urbana vs rural, climatica, mix de clientes).
Esto degrada la convergencia. En este trabajo se aborda parcialmente
reduciendo el numero de rondas locales y monitoreando el gap contra
el modelo centralizado (capitulo 5).

\subsection{Regulacion electrica chilena y los cargos SEC/SERNAC}

La regulacion electrica chilena se estructura en torno a:

\begin{itemize}[leftmargin=*,itemsep=2pt]
  \item \textbf{CEN} (Coordinador Electrico Nacional): opera el SEN,
        coordina despacho y publica medidas de calidad.
  \item \textbf{CNE} (Comision Nacional de Energia): regula normas
        tecnicas y tarifarias.
  \item \textbf{SEC} (Superintendencia de Electricidad y Combustibles):
        fiscaliza cumplimiento de calidad, seguridad y normativa
        tecnica. Emite cargos directos por incumplimiento.
  \item \textbf{SERNAC}: protege derechos del consumidor. Coordina
        con la SEC los casos de reposicion tardia o dano a usuarios.
\end{itemize}

Los indicadores \emph{SAIDI} (System Average Interruption Duration
Index) y \emph{SAIFI} (System Average Interruption Frequency Index)
son los principales metricas de calidad. La regulacion del NTIC
(Norma Tecnica de Calidad de Servicio) fija valores maximos por zona
y compensacion automatica a clientes en caso de superarlos.

Tras el frente de Valparaiso (julio 2026), la SEC y el SERNAC
comenzaron a emitir cargos y procedimientos contra Chilquinta y otras
concesionarias por \emph{reposicion tardia}. El plazo maximo de
reposicion en zona urbana es 8~h; en zona rural 24~h. Cualquier
sistema que ayude a reducir el tiempo medio de reposicion reduce
directamente la exposicion regulatoria.

\subsection{Trabajo relacionado}

\subsubsection{Detectores no supervisados en subestaciones}

La mayoria de la literatura usa \emph{autoencoders variacionales} (VAE)
o \emph{normalizing flows} sobre los datos crudos SCADA~\citep{AE-IEA-2021}.
El aporte de este trabajo es que el residuo sobre el gemelo digital es
\emph{mas robusto al punto de operacion} que la senal cruda, porque
los cambios legitimos de consigna son predichos por el gemelo y se
cancelan en el residuo.

\subsubsection{Conformal prediction para series de tiempo}

La \emph{conformal prediction adaptativa}~\citep{ACI-2020} actualiza
los cuantiles con una ventana deslizante. En este trabajo se usa la
version estatica (cuantil fijo sobre calibracion inicial) por
simplicidad y reproducibilidad, dejando la version adaptativa como
trabajo futuro.

\subsubsection{Aprendizaje federado en sistemas electricos}

La literatura es aun escasa. \citet{FedGrid-2022} propone FedAvg
para deteccion de robo de energia en smart meters. Este trabajo
extiende el enfoque a subestaciones, donde la dimensionalidad y la
criticidad de las fallas es muy distinta.

\subsection{Sintesis}

El marco teorico presentado cubre los cuatro pilares tecnologicos del
trabajo (gemelo digital, conformal prediction, federated learning) y los
fundamentos regulatorios que dan sentido economico al producto. El
capitulo siguiente presenta los datos concretos sobre los que se
valida el pipeline.